\subsection{QUOTIENT RINGS}

Let A be a ring. If a is a two-sided ideal of A, two elements x and y of A are said to be congruent modulo a, written \(x \equiv y \pmod{a}\) or \(x \equiv y(a)\) , if \(x - y \in a\) . This is an equivalence relation on A. The relations \(x \equiv y(a)\) and \(x' \equiv y'(a)\) imply \(x + x' \equiv y + y'(a)\) , \(xx' \equiv xy'(a)\) for a is a left ideal and \(xy' \equiv yy'(a)\) for a is a right ideal, whence \(xx' \equiv yy'(a)\) . Conversely, if R is an equivalence relation on A compatible with addition and multiplication, the set a of \(x \equiv 0 \mod R\) is a two-sided ideal and \(x \equiv y \mod R\) is equivalent to \(x \equiv y \mod a\) .

Let A be a ring and a two-sided ideal of A. A/a denotes the quotient set of A by the equivalence relation \(x \equiv y(a)\) , with addition and multiplication the quotients of those on A (§ 1, no. 6, Definition 11). We show that A/a is a ring:

\begin{enumerate}
\def\labelenumi{(\alph{enumi})}
\item
  Under addition, A/a is the quotient commutative group of the additive group of A by the subgroup a.
\item
  Under multiplication, A/a is a monoid (§ 2, no. 1).
\item
  Let \(\xi, \eta, \zeta\) be in A/a and let \(\pi: \mathbf{A} \to \mathbf{A}/\mathfrak{a}\) be the canonical mapping; we choose elements \(x, y, z\) in A such that \(\pi(x) = \xi, \pi(y) = \eta\) and \(\pi(z) = \zeta\) . Then
\end{enumerate}

\[
\begin{array}{r l} \xi (\eta + \zeta) & = \pi (x) \pi (y + z) = \pi (x (y + z)) = \pi (x y + x z) \\ & = \pi (x) \pi (y) + \pi (x) \pi (z) = \xi \eta + \xi \zeta \end{array}
\]

and the relation \((\xi + \eta)\zeta = \xi \zeta + \eta \zeta\) is established similarly.

DEFINITION 6. Let A be a ring and a a two-sided ideal of A. The quotient ring of A by a, denoted by A/a, is the quotient set of A by the equivalence relation \(x \equiv y(a)\) , with addition and multiplication the quotients of those on A.

The ring A/\{0\} is isomorphic to A and A/A is a zero ring.

THEOREM 2. Let A be a ring and a two-sided ideal of A.

\begin{enumerate}
\def\labelenumi{(\alph{enumi})}
\item
  The canonical mapping \(\pi\) of A onto A/a is a ring homomorphism.
\item
  Let \(\mathbf{B}\) be a ring and \(f\) a homomorphism of \(\mathbf{A}\) into \(\mathbf{B}\) . If \(f(\mathfrak{a}) = \{0\}\) , there exists one and only one homomorphism \(\bar{f}\) of \(\mathbf{A} / \mathfrak{a}\) into \(\mathbf{B}\) such that \(f = \bar{f} \circ \pi\) .
\end{enumerate}

By construction, \(\pi(x + y) = \pi(x) + \pi(y)\) and \(\pi(xy) = \pi(x)\pi(y)\) for \(x, y\) in A; also \(\pi(1)\) is the unit \(\varepsilon\) of A/a, whence (a).

Let \(\mathbf{A}^{+}\) be the additive group of \(\mathbf{A}\) and \(\mathbf{B}^{+}\) that of \(\mathbf{B}\) ; as \(f\) is a homomorphism of \(\mathbf{A}^{+}\) into \(\mathbf{B}^{+}\) , zero on the subgroup \(a\) of \(\mathbf{A}^{+}\) , there exists (§ 4, no. 4, Proposition 5) one and only one homomorphism \(\bar{f}\) of \(\mathbf{A}^{+}/a\) into \(\mathbf{B}^{+}\) such that \(f = \bar{f} \circ \pi\) . Let \(\xi, \eta\) be in \(\mathbf{A}/a\) ; choose \(x, y\) in \(\mathbf{A}\) with \(\pi(x) = \xi\) and \(\pi(y) = \eta\) ; then \(\xi \eta = \pi(xy)\) , whence

\[
\bar {f} (\xi \eta) = \bar {f} (\pi (x y)) = f (x y) = f (x). f (y) = \bar {f} (\xi). \bar {f} (\eta)
\]

and \(\bar{f}(\varepsilon) = f(\pi(1)) = f(1)\) , hence \(\bar{f}\) is a ring homomorphism.

THEOREM 3. Let A and B be rings and f a homomorphism of A into B.

\begin{enumerate}
\def\labelenumi{(\alph{enumi})}
\item
  The kernel \(\mathfrak{a}\) off is a two-sided ideal of \(\mathbf{B}\) .
\item
  The image \(\mathbf{B}' = f(\mathbf{B})\) off is a subring of \(\mathbf{B}\) .
\item
  Let \(\pi: \mathbf{A} \to \mathbf{A} / \mathfrak{a}\) and \(i: \mathbf{B}' \to \mathbf{B}\) be the canonical morphisms. There exists one and only one morphism \(\bar{f}\) of \(\mathbf{A} / \mathfrak{a}\) into \(\mathbf{B}'\) such that \(f = i_{\circ} \bar{f} \circ \pi\) and \(\bar{f}\) is an isomorphism.
\end{enumerate}

As \(f\) is a morphism of the additive group of \(\mathbf{A}\) into that of \(\mathbf{B}\) , \(\mathfrak{a}\) is a subgroup of \(\mathbf{A}\) . If \(x \in \mathfrak{a}\) and \(a \in \mathbf{A}\) , then \(f(ax) = f(a)f(x) = 0\) , hence \(ax \in \mathfrak{a}\) and similarly \(xa \in \mathfrak{a}\) ; hence \(\mathfrak{a}\) is a two-sided ideal of \(\mathbf{A}\) . Assertion (b) is obvious. As \(f\) is zero on \(\mathfrak{a}\) , there exists a morphism \(\bar{f}\) of \(\mathbf{A} / \mathfrak{a}\) into \(\mathbf{B}'\) such that \(f = i \circ \bar{f} \circ \pi\) (Theorem 2). The uniqueness of \(\bar{f}\) and the fact that \(\bar{f}\) is an isomorphism follow from Set Theory, II, § 6, no. 4.

